\chapter{Literature Review}\label{chap:litreview}

This chapter reviews the literature thematically. Its argument is: secure neural-network training under \ac{mpc} is feasible but costly; the cost depends on the operations, the protocol and the system; \ac{ff} changes the computation graph of learning and has been studied and optimised in non-cryptographic settings; the privacy work on \acp{snn} identified in this review is adjacent (encrypted evaluation or training-data leakage) rather than a secure-\ac{mpc} comparison; and a controlled experimental comparison of \ac{ff} and \ac{bp} under one specified \ac{mpc} configuration is therefore worth testing. This is a \emph{preliminary} gap derived from a non-exhaustive search, not a claim of novelty.

\paragraph{Scope and sources.} The review covers the 23 references of the approved project proposal, all of which were read by the author, and three later additions that are distinguished as supplementary: two works from 2026 (Zhou \cite{zhou2026rmfe}; Elnady et al.\ \cite{elnady2026ff}) and the multi-party training framework of Lu, Basaran and G\"uler \cite{lu2024clover}, whose introduction, problem formulation and experimental section were read for this version. One further entry supports a specific technical statement rather than the thematic review: the probabilistic-truncation analysis of Santos et al.\ \cite{santos2024curl}, characterised from its abstract. The full-text files of six proposal references \cite{keller2023note,luo2026secptuning,yan2025comet,feng2025fwdp,nikfam2025spyking,kim2022privatesnn} are not part of the project archive. The characterisations of these six therefore follow the project's literature records and the official abstracts. Numerical results quoted below always carry the setting of the paper that reports them; they use different protocols, hardware, networks, models and baselines and are not comparable with each other or with the measurements of this thesis.

\section{MPC for Privacy-Preserving Machine Learning}\label{sec:lr-mpc}
The survey by Zhou et al.\ organises secure machine learning by security model, system type and deployment model and discusses secret sharing and homomorphic encryption. In the systems it surveys, secret-sharing approaches often compute faster than homomorphic-encryption approaches but may need more communication; this is a trend in the reviewed work and not a universal law \cite{zhou2024smpc}. It is a survey and does not benchmark \ac{ff}.

Two frameworks matter for the experimental part. MP-SPDZ provides a high-level interface, a compiler and a virtual machine for many \ac{mpc} protocols and security settings, so that a computation can be written once and run under different assumptions \cite{keller2020mpspdz}; it is the implementation framework of this project, and statements about security always have to name the selected protocol. CrypTen offers a PyTorch-like tensor and automatic-differentiation interface under a semi-honest model in which a trusted party can supply correlated randomness; its published evaluations concentrate on private inference rather than on a comprehensive secure-training study \cite{knott2021crypten}.

Neighbouring paradigms must not be confused with \ac{mpc}. Salam et al.\ combine homomorphic encryption with federated collaboration \cite{salam2023mpppml}; this is a different mechanism from secret-sharing \ac{mpc}, and it is mentioned only to illustrate the paradigm. Federated learning on its own does not give a cryptographic guarantee, and \ac{dp} protects released outputs, not inputs during computation.

\section{Secure Training Systems and Baselines}\label{sec:lr-systems}
Keller and Sun build quantised neural-network training on MP-SPDZ and contribute secure protocols for operations such as exponentiation and inverse square root. In their three-party honest-majority experiment a convolutional MNIST network reaches about 99.2\% accuracy in about 3.5 hours; further networks and datasets are also evaluated, and the paper discusses numerical error and the cost of nonlinear operations \cite{keller2022sqt}. A later note by the same authors identifies an implementation error in the max-pooling backward-gradient routing that explains an accuracy gap observed in the earlier work \cite{keller2023note}; it is a correction note and not a claim that quantisation never changes accuracy. This framework is the closest methodological precedent for the present work, although its models are far larger than the bounded prototype used here.

Falcon is a three-party honest-majority framework with semi-honest and malicious (with abort) variants. In its own private-training setup it reports about $6\times$ faster training than SecureNN and $4.4\times$ than ABY3 with $2$--$60\times$ less communication, finds that compute can dominate for large networks in a wide-area setting, and excludes inference attacks on released outputs \cite{wagh2021falcon}. CryptGPU runs secure learning on GPUs in a replicated-sharing setting and reports speedups against specific prior secure baselines; some comparisons depend on batch size and extrapolated training settings \cite{tan2021cryptgpu}. ESMPC targets secure training in two- and three-party semi-honest settings, profiles bottlenecks under the Piranha framework and introduces pipeline and protocol optimisations; against its Piranha-based baseline its abstract reports training-time reductions of over 30\% for AlexNet and over 45\% for FCNN, LeNet and ViT in the two-party setting, and over 45\% for all four models in the three-party setting \cite{yang2025esmpc}. SWIFT provides maliciously secure three-party computation over rings with \ac{god} in an honest-majority setting, plus a four-party extension. It benchmarks logistic-regression training and inference and deep-network \emph{inference} (LeNet, VGG16) in a 64-bit-ring wide-area setting; neural-network training would require mixed-world conversions that the paper leaves as future work \cite{koti2021swift}. The more recent work of Zhou uses Shamir sharing over Galois rings with reverse multiplication-friendly embedding (RMFE) packing and mixed-circuit techniques for comparisons, truncation, activations, pooling and loss computation. It trains fully connected and convolutional networks on MNIST for at least three parties with an honest majority of passively corrupted parties; relative to its Galois-ring packing baseline at equal packing density it reports $71.4$--$75.1\%$ lower communication and $55.3$--$70.1\%$ faster online training in a local-area network, and $5.3$--$68.1\%$ faster in a wide-area network across its three networks, on its own hardware and network settings \cite{zhou2026rmfe}. It shows that secure-training optimisation is still advancing but it does not study \ac{ff}, and its protocol differs from the replicated-ring configuration used here.

Scalability in the number of parties is a separate axis. Lu, Basaran and G\"uler argue that information-theoretic multi-party training frameworks based on Lagrange coded computing suffer from communication that grows quadratically with the number of users, because polynomial degree reduction is needed after every multiplication. Their framework, CLOVER, introduces a degree-reduction mechanism (Double Lagrange Coding) that moves the quadratic part of the communication to a data-independent offline phase and achieves linear communication complexity under information-theoretic privacy against up to $T$ colluding honest-but-curious users \cite{lu2024clover}. It operates over a prime field and, to remain polynomial, trains neural networks with quadratic activations rather than ReLU; its experiments train a one-hidden-layer network with 128 hidden units on MNIST and on VGG features of CIFAR-10 and report, against its COPML baseline, about $28\times$ and $26\times$ lower communication and $6\times$ and $4.8\times$ faster training, respectively, in its own distributed setting \cite{lu2024clover}. The work is relevant here because it treats communication as the scaling bottleneck of multi-party training and because it separates offline from online communication, but it changes the network (polynomial activations) and the sharing scheme, and it does not study \ac{ff}.

Table~\ref{tab:lr-mpc} summarises these works with the setting and the limits that matter for this thesis.

\begin{table}[ht]
\centering
\footnotesize
\begin{tabular}{@{}>{\raggedright\arraybackslash}p{2.2cm}>{\raggedright\arraybackslash}p{3.8cm}>{\raggedright\arraybackslash}p{3.6cm}>{\raggedright\arraybackslash}p{4.0cm}@{}}
\toprule
Work & Setting and security model & Task & Limits for this thesis \\
\midrule
Keller--Sun \cite{keller2022sqt} & MP-SPDZ; three parties, honest majority (main experiment) & Quantised training; MNIST CNN and others & Models far larger than the bounded prototype; not \ac{ff} \\
Falcon \cite{wagh2021falcon} & Three parties; semi-honest and malicious with abort & Private training and inference & Outputs not protected; own network setting \\
CryptGPU \cite{tan2021cryptgpu} & Replicated sharing, GPU & Secure learning & Baseline- and batch-size-dependent speedups \\
ESMPC \cite{yang2025esmpc} & Two and three parties, semi-honest, GPU & Training: FCNN, LeNet, AlexNet, ViT & Gains relative to its Piranha baseline \\
SWIFT \cite{koti2021swift} & Three/four parties, honest majority, malicious, \acs{god} & Logistic-regression training and inference; DNN inference & DNN training left as future work \\
Zhou \cite{zhou2026rmfe} & Shamir over Galois rings, $n\ge3$, semi-honest, $t<n/2$ & Training: fully connected and convolutional MNIST networks & Different sharing scheme and environment; not \ac{ff} \\
CrypTen \cite{knott2021crypten} & Semi-honest; trusted party for randomness & Mainly private inference & Not a comprehensive training study \\
CLOVER \cite{lu2024clover} & Lagrange coded computing over a prime field; $N$ users, up to $T$ honest-but-curious colluders & Training of polynomial (quadratic-activation) networks; MNIST, CIFAR-10 features & Different network and sharing scheme; not \ac{ff} \\
\bottomrule
\end{tabular}
\caption{Secure-training and inference systems relevant to this thesis, as characterised in the cited papers.}
\label{tab:lr-mpc}
\end{table}

\section{Protocol-Level and System Optimisations}\label{sec:lr-opt}
Optimisation work has been carried out at several layers of the stack while keeping the learning algorithm fixed. At the system layer, GPU acceleration \cite{tan2021cryptgpu} and compiler or runtime design \cite{smajlovic2023sequre} reduce time; Sequre targets biomedical secure-computation workloads and relies on an auxiliary trusted dealer for correlated randomness, an additional assumption that must be stated whenever its results are discussed \cite{smajlovic2023sequre}. At the protocol layer, Zhang et al.\ use Shamir sharing and propose secure fixed-point arithmetic, a Winograd-based convolution that reduces the multiplications of one $2\times2$-output, $3\times3$-kernel convolution from 36 to 16, and an approximate softmax; their evaluation concerns protocols and primitives rather than complete training runs \cite{zhang2025efficient}.

Das et al.\ show that the learning procedure and the protocol interact: their two-party protocol combines \ac{mpc} with \ac{dp} and reports $56$--$794\times$ lower communication and $16$--$182\times$ lower runtime than earlier secure-and-private protocols in its selected comparisons, where some baseline runtimes are lower-bound estimates because earlier code was unavailable \cite{das2025dpmpc}. This is a \ac{dp}-\ac{mpc} result and not a speedup of pure \ac{mpc}.

A common observation is that the cost of a secure operation is highly non-uniform and that the learning algorithm determines which operations occur. SecP-Tuning is \ac{mpc}-based forward-only prompt tuning of \acp{llm} with a data-owner/server interaction; it reduces the cost of backward propagation, optimisation and softmax attention in its setting, and its abstract reports roughly $12\times$ and $16\times$ end-to-end acceleration and $17\times$ and $20\times$ lower communication against full-parameter fine-tuning and gradient-based prompt tuning, respectively \cite{luo2026secptuning}. It is not the \ac{ff} algorithm, and its offloading and trust model differ from this thesis. Comet is \ac{mpc}-based private \emph{inference} for \acp{llm}: it predicts activation sparsity, protects the input-dependent sparsity pattern, and reports $1.87$--$2.63\times$ speedup and $1.94$--$2.64\times$ lower communication in its workloads \cite{yan2025comet}. Neither paper evaluates secure training with \ac{ff}.

\section{Fixed-Point and Floating-Point Training}\label{sec:lr-numeric}
The numerical representation affects both accuracy and cost. Keller and Sun develop fixed-point training with probabilistic truncation \cite{keller2022sqt}; CrypTen and Zhang et al.\ discuss overflow, wrap-around and truncation effects \cite{knott2021crypten,zhang2025efficient}. BEACON provides secure floating-point training in a two-party semi-honest setting and reports multi-fold improvements over the floating-point libraries and protocols it compares against; its numbers cannot be read as a general advantage over fixed-point training or over other frameworks \cite{rathee2023beacon}. For \ac{ff}, which uses squared activities, normalisation and a nonlinear loss, ranges and precision are particular concerns; no source reviewed here evaluates them under \ac{mpc}.

\section{Forward-Forward and Local Learning}\label{sec:lr-ff}
Hinton introduces \ac{ff} with positive and negative passes and layer-local goodness objectives. The paper evaluates relatively small networks, finds \ac{ff} competitive in some settings and worse or slower than \ac{bp} in others, and describes the study as preliminary, leaving larger-scale behaviour open \cite{hinton2022ff}. Self-Contrastive \ac{ff} generates positive and negative examples from the data themselves for unsupervised local learning and reports, in its settings, $98.70\pm0.01\%$ on MNIST with a two-hidden-layer perceptron (2{,}000 units per layer), $99.37\pm0.06\%$ with a convolutional network, $80.75\pm0.12\%$ on CIFAR-10 and $77.30\pm0.12\%$ on STL-10; on STL-10 an additional fine-tuning stage uses end-to-end backpropagation, so the extended pipeline is not purely local \cite{chen2025scff}.

Elnady et al.\ compare \ac{ff} with \ac{bp} on six datasets and tune \ac{ff} hyperparameters with two conventional and four metaheuristic methods on four of them, reporting validation-accuracy gains (for example 93\% to 97\% on Pneumonia-MedMNIST and 87\% to 91\% on Fashion-MNIST) together with comparable or reduced training time in their settings \cite{elnady2026ff}. Their \ac{ff} runs use Hinton's implementation while \ac{bp} uses the default training function of a numerical-computing toolbox, so the two methods are not necessarily matched in implementation or computational boundary, and the work involves no cryptography.

Forward-learning methods have also been combined with differential privacy. The method of Feng et al.\ adds noise and privacy guarantees to forward learning, but it is not the \ac{ff} algorithm and it does not use \ac{mpc} \cite{feng2025fwdp}.

\section{Distributed and Federated Forward-Forward}\label{sec:lr-dist}
The layer-wise structure of \ac{ff} has been exploited for parallelism. \ac{pff} distributes layer training over four compute nodes; in one MNIST experiment an all-layers configuration with adaptive negative data reaches $98.51\%$ against $98.52\%$ for the sequential baseline, in $2980.76$\,s against $11190.72$\,s, a $3.75\times$ speedup. It uses no cryptography, and its faster softmax variants add a classifier trained with backpropagation \cite{aktemur2025pff}. A pipeline-parallel speedup does not establish a speedup under \ac{mpc}, because it concerns waiting time between compute nodes and not secure-operation cost. FedFwd integrates \ac{ff}-style local training into federated learning and compares it with backpropagation-based federated methods; its results depend on the data distribution (i.i.d.\ and non-i.i.d.), and it is a workshop paper \cite{park2023fedfwd}. Federated learning alone is not a cryptographic guarantee and is not \ac{mpc}.

\section{Spiking Neural Networks and Privacy}\label{sec:lr-snn}
Aribe reviews surrogate-gradient training, \ac{ann}-to-\ac{snn} conversion and \ac{stdp} and notes that some hardware figures are taken from other literature rather than measured by the author \cite{aribe2025snn}; such figures are not used as evidence of cost under \ac{mpc}. SpyKing uses BFV fully homomorphic encryption, not \ac{mpc}: LeNet5 and a spiking counterpart are trained in a PyTorch phase and then evaluated on encrypted data, and ReLU and other nonlinear operations require decrypting and re-encrypting intermediate values in the described implementation, so the pipeline is not end-to-end encrypted through its nonlinear layers; the cost of the spiking model depends on the temporal sequence length \cite{nikfam2025spyking}. PrivateSNN addresses leakage of training data during conversion from an \ac{ann} and class leakage from the converted parameters, using synthetic samples; within its own setting it reports a performance drop below about 2\% and a large energy-efficiency gain, and it is not a secure-computation protocol \cite{kim2022privatesnn}. \ac{fhe} evaluation of \acp{snn} does not establish \ac{snn} training under \ac{mpc}, and biological inspiration is neither a privacy guarantee nor, on this evidence, automatically cheaper under encryption. This topic is optional for the thesis.

\section{Work at the Intersection of MPC and Biologically Plausible Learning}\label{sec:lr-intersection}
Adjacent work touches the intersection from several sides: SecP-Tuning removes backward computation in \ac{mpc}-based prompt tuning \cite{luo2026secptuning}; Comet exploits activation sparsity for \ac{mpc} inference \cite{yan2025comet}; \ac{ff} has been parallelised and used in federated training without cryptography \cite{aktemur2025pff,park2023fedfwd} and optimised without cryptography \cite{elnady2026ff}; and \acp{snn} have been studied under \ac{fhe} \cite{nikfam2025spyking}. Table~\ref{tab:lr-adjacent} places these side by side.

\begin{table}[ht]
\centering
\footnotesize
\begin{tabular}{@{}>{\raggedright\arraybackslash}p{2.5cm}>{\raggedright\arraybackslash}p{2.8cm}>{\raggedright\arraybackslash}p{2.6cm}>{\raggedright\arraybackslash}p{2.9cm}>{\raggedright\arraybackslash}p{2.0cm}@{}}
\toprule
Work & Learning rule & Privacy mechanism & Task & \ac{mpc} training of \ac{ff}? \\
\midrule
Keller--Sun \cite{keller2022sqt} & \ac{bp} & \ac{mpc} & Training & No \\
Zhou \cite{zhou2026rmfe} & \ac{bp}-style & \ac{mpc} (Shamir) & Training & No \\
SecP-Tuning \cite{luo2026secptuning} & Forward-only, prompt tuning & \ac{mpc} & \ac{llm} fine-tuning & No \\
Comet \cite{yan2025comet} & -- & \ac{mpc} & \ac{llm} inference & No \\
\ac{pff} \cite{aktemur2025pff} & \ac{ff} & None & Distributed training & No \\
FedFwd \cite{park2023fedfwd} & \ac{ff} & Federated (no cryptography) & Federated training & No \\
Elnady et al.\ \cite{elnady2026ff} & \ac{ff} vs.\ \ac{bp} & None & Tuning, six datasets & No \\
SpyKing \cite{nikfam2025spyking} & \ac{snn}/\ac{ann}, trained in plaintext & \ac{fhe} (BFV) & Encrypted evaluation & No \\
\bottomrule
\end{tabular}
\caption{Adjacent work and its distance from \ac{mpc}-based \ac{ff} training, as characterised in the cited papers.}
\label{tab:lr-adjacent}
\end{table}

\section{Comparative Synthesis and Research Gap}\label{sec:lr-gap}
The \ac{mpc} literature improves the cryptographic and systems stack, and recent work such as that of Zhou shows that this is an active area; the \ac{ff} literature studies and optimises the learning algorithm in non-cryptographic settings, for example in the work of Elnady et al. Within the works reviewed for this thesis (proposal literature search up to 11 October 2026, using general web search with the terms Forward-Forward, local or layer-wise learning, spiking neural networks, gating or sparsity, secure multi-party computation, secret sharing and private training), the reviewed literature did not identify a controlled, end-to-end, matched comparison of \ac{ff} and \ac{bp} under the \emph{same specified} \ac{mpc} framework, protocol and threat model, in which communication, runtime and utility are evaluated together, followed by a bottleneck-driven \ac{mpc}-aware modification of \ac{ff}. The supplementary work added later (Zhou, Elnady et al., Lu et al.) does not change this observation: each improves either the secure-training stack or the plaintext learning algorithm, but none evaluates \ac{ff} under \ac{mpc}. This is a preliminary observation: the systematic search has not yet been repeated after 11 October 2026, the search was not exhaustive, and closely related unpublished or recent work may exist; the search must be refreshed and its date updated before submission. It motivates the experimental study of research questions 1--3 but is not a claim that no related work exists.
